\section{问题四：离网运行与储能配置研究}

\subsection{模型建立}

问题四将园区由并网运行转为离网运行，此时外部电网不再补足功率缺额，
园区用能完全受风光发电功率约束。与问题三相同，电氢氨装置运行时功率
连续可调且最低稳定负荷为 10\%；但在离网且风光不足时，装置不能依靠电网
维持最低负荷，因此允许停机。设 $\alpha_{t,s}$ 为场景 $s$ 下第 $t$ 小时
制氨负荷率，则
\begin{equation}
  \alpha_{t,s}=0
  \quad\text{或}\quad
  0.1\leq\alpha_{t,s}\leq1.
  \label{eq:p4-alpha}
\end{equation}
式~\eqref{eq:p4-alpha} 表示：若装置运行，则不得低于 10\% 负荷；
若风光出力连最低稳定负荷也无法支撑，则该时段停机。

无储能时，风光首先满足园区常规负荷，剩余功率用于制氨。记扩容后制氨
系统满负荷用电功率为
\begin{equation}
  P^{\max}=20+20+1.5=41.5~\mathrm{MW}.
  \label{eq:p4-pmax}
\end{equation}
在尽可能利用风光发电的条件下，负荷率按下式确定：
\begin{equation}
  \alpha_{t,s}=
  \begin{cases}
    0, &
    P_{t,s}^{\mathrm{re}}-P_{t,s}^{\mathrm{load}} < 0.1P^{\max},\\[2mm]
    \min\left\{1,\dfrac{P_{t,s}^{\mathrm{re}}-P_{t,s}^{\mathrm{load}}}{P^{\max}}\right\}, &
    P_{t,s}^{\mathrm{re}}-P_{t,s}^{\mathrm{load}} \geq 0.1P^{\max}.
  \end{cases}
  \label{eq:p4-no-storage-alpha}
\end{equation}
对应弃电功率和常规负荷缺额为
\begin{align}
  P_{t,s}^{\mathrm{curtail}}
  &=\max\left\{P_{t,s}^{\mathrm{re}}-P_{t,s}^{\mathrm{load}}
  -P^{\max}\alpha_{t,s},0\right\},\\
  P_{t,s}^{\mathrm{unserved}}
  &=\max\left\{P_{t,s}^{\mathrm{load}}-P_{t,s}^{\mathrm{re}},0\right\}.
  \label{eq:p4-curtail-unserved}
\end{align}
日产氨量为
\begin{equation}
  Q_s=\sum_{t=1}^{24}3\alpha_{t,s}.
  \label{eq:p4-production}
\end{equation}

为量化园区能源自治所需的最小可再生能源装机容量，本文将能源自治界定为基础用电保障能力，
即在全时段离网运行场景下，风光发电出力始终不低于园区常规负荷需求。在此约束框架下，
若风电与光伏装机容量按统一扩容系数 $k$配置，则
\begin{equation}
  k\left(P_{t,s}^{\mathrm{wind}}+P_{t,s}^{\mathrm{pv}}\right)
  \geq P_{t,s}^{\mathrm{load}},\qquad \forall s,t.
  \label{eq:p4-autonomy}
\end{equation}
由此可得最小等比例扩容倍数
\begin{equation}
  k_{\min}=
  \max_{s,t}
  \frac{P_{t,s}^{\mathrm{load}}}{P_{t,s}^{\mathrm{re}}}.
  \label{eq:p4-kmin}
\end{equation}
本模型将“能源自治”界定为保障园区常规负荷的持续供电，不要求全时段满负荷制氨；
否则会使装机规模过度放大，偏离“最小自治容量”的估算目标。

储能参与时，针对无储能结果中的最大弃电场景配置储能。引入充电功率
$P_t^{\mathrm{ch}}$、放电功率 $P_t^{\mathrm{dis}}$、储能电量
$E_t^{\mathrm{ess}}$、储能容量 $E^{\max}$ 和启停变量 $y_t$，有
\begin{equation}
  0.1y_t\leq\alpha_t\leq y_t,\qquad y_t\in\{0,1\}.
  \label{eq:p4-run-status}
\end{equation}
离网功率平衡为
\begin{equation}
  P_t^{\mathrm{re}}+P_t^{\mathrm{dis}}+P_t^{\mathrm{unserved}}
  =
  P_t^{\mathrm{load}}+P^{\max}\alpha_t+P_t^{\mathrm{ch}}+P_t^{\mathrm{curtail}},
  \label{eq:p4-storage-balance}
\end{equation}
储能状态方程为
\begin{equation}
  E_{t+1}^{\mathrm{ess}}
  =(1-\lambda)E_t^{\mathrm{ess}}
  +\eta_cP_t^{\mathrm{ch}}\Delta t
  -\frac{P_t^{\mathrm{dis}}\Delta t}{\eta_d},
  \qquad
  0\leq E_t^{\mathrm{ess}}\leq E^{\max}.
  \label{eq:p4-storage-state}
\end{equation}
其中 $\eta_c=\eta_d=0.9$，$\lambda=0.2\%$。储能配置不单独追求最大消纳，
而是在候选功率--容量组合中同时比较经济折中方案和消纳导向方案：前者以吨氨
成本改善和较小投资为主，后者以最大弃电场景弃电量最小为主。储能投资成本按
使用寿命折算到日成本：
\begin{equation}
  C_{\mathrm{ESS}}^{\mathrm{day}}
  =
  \frac{1000c_{\mathrm{ESS}}^{\mathrm{cap}}E^{\max}}{15\times365}
  +1000c_{\mathrm{ESS}}\sum_tP_t^{\mathrm{dis}}\Delta t,
  \label{eq:p4-storage-cost}
\end{equation}
其中 $E^{\max}$ 单位为 MWh，$c_{\mathrm{ESS}}^{\mathrm{cap}}=1000$ 元/kWh，
储能运维系数 $c_{\mathrm{ESS}}=0.01$ 元/kWh。储能成本按功率、容量和寿命折算
是储能技术经济分析中的常用处理方式~[4]。吨氨成本同时计入风光
度电成本、设备运维、合成氨固定投资折算成本以及储能年化成本。

\subsection{模型求解与结果}

首先计算无储能离网运行结果。原始风光装机下，24 个场景全年折算制氨总量为
9781.56~t，年平均产能利用率为 $37.74\%$，全年弃电量为 14188.67~MWh。
由于个别低出力时段风光甚至不能覆盖常规负荷，全年常规负荷缺额为
8.29~MWh。由式~\eqref{eq:p4-kmin} 得到仅覆盖园区常规负荷的风光等比例
扩容倍数为 $k_{\min}=1.0926$，对应风电装机约为 43.71~MW，光伏装机约为
69.93~MW。若不强制保持风光装机比例，而是在当前 40~MW 风电、64~MW 光伏
基础上最小化新增容量，网格搜索得到的最小自治配置为风电 43.8~MW、光伏
64.0~MW，即仅需新增约 3.8~MW 风电即可覆盖常规负荷。

\begin{figure}[H]
  \centering
  \includegraphics[width=0.88\textwidth]{p4_offgrid_dispatch_example.pdf}
  \caption{无储能离网条件下代表场景的风光利用与制氨负荷率}
  \label{fig:p4-offgrid-dispatch}
\end{figure}

图~\ref{fig:p4-offgrid-dispatch} 展示了最大弃电场景 $W4\_P1$ 的无储能调度。
当风光出力高于常规负荷和制氨满负荷需求时，制氨负荷率达到 1，剩余功率被弃电；
当风光出力不足以支撑最低稳定负荷时，装置停机。这是离网条件下产能利用率
低于并网情形的重要原因。

无储能结果中，最大弃电场景为 $W4\_P1$。针对该场景进行储能候选配置搜索，
经济折中方案选择 \textbf{5~MW/5~MWh} 储能。该方案下，$W4\_P1$ 场景日产氨量
由 46.19~t 提高至 46.77~t，日弃电量由 177.29~MWh 降至 164.30~MWh，
吨氨成本由 4256.93~元/t 降至 4248.95~元/t。若以最大消纳为目标，
40~MW/160~MWh 储能可将 $W4\_P1$ 弃电量降为 0，并将日产氨量提高至
56.41~t，但吨氨成本升至 4337.38~元/t。因此本文以经济折中方案作为
推荐储能配置。储能充放电和能量状态如图~\ref{fig:p4-storage-dispatch} 所示。

\begin{figure}[H]
  \centering
  \includegraphics[width=0.88\textwidth]{p4_storage_dispatch.pdf}
  \caption{最大弃电场景下储能参与后的充放电与荷电状态}
  \label{fig:p4-storage-dispatch}
\end{figure}

将该小容量储能用于 24 个场景重新调度，全年折算制氨总量提高至
9969.43~t，年平均产能利用率提高至 $38.46\%$，常规负荷缺额降为 0，
弃电量由 14188.67~MWh 降至 10983.99~MWh。各场景弃电量对比如
图~\ref{fig:p4-curtailment} 所示。结果表明，小容量储能主要承担削减短时
缺额和部分转移富余电量的作用，能够改善能源自治与经济性，但无法完全吸收
所有高峰风光出力。

\begin{figure}[H]
  \centering
  \includegraphics[width=0.9\textwidth]{p4_curtailment_by_scenario.pdf}
  \caption{无储能与有储能离网运行下各场景日弃电量对比}
  \label{fig:p4-curtailment}
\end{figure}

表~\ref{tab:p4-mode-compare} 比较了联网连续调节、离网无储能和离网有储能三种
运行方式。联网连续调节复用问题三成本最低方案结果。离网无储能的吨氨成本较低，
但存在常规负荷缺额，且全年制氨量明显低于联网方案；配置小容量储能后，常规负荷
缺额被消除，制氨量和吨氨成本均有所改善，但弃电仍然存在。

\begin{table}[H]
  \centering
  \caption{联网与离网运行方式的全年结果对比}
  \label{tab:p4-mode-compare}
  \begin{tabular}{lccc}
    \toprule
    运行方式 & 全年制氨量(t) & 吨氨成本(元/t) & 年产能利用率 \\
    \midrule
    联网连续调节 & 12960.00 & 4302.07 & 50.00\% \\
    离网无储能 & 9781.56 & 4178.99 & 37.74\% \\
    离网有储能 & 9969.43 & 4167.49 & 38.46\% \\
    \bottomrule
  \end{tabular}
\end{table}

\begin{figure}[H]
  \centering
  \includegraphics[width=0.78\textwidth]{p4_mode_cost_compare.pdf}
  \caption{联网与离网运行方式的吨氨成本和全年制氨量对比}
  \label{fig:p4-mode-compare}
\end{figure}

综合来看，原始风光装机下离网运行可以实现部分能源自治，但受低出力时段制约，
制氨产量和产能利用率下降。小容量储能能够消除常规负荷缺额、减少弃电并略微
降低吨氨成本；大容量储能则可进一步提高风光消纳，但投资成本会削弱经济性。
因此，联网方案在产量稳定性上仍具有优势，而离网储能方案的作用主要体现在提升
能源自治能力和局部风光利用水平。
